\documentclass[letterpaper]{article} 
\usepackage[utf8]{inputenc}
\linespread{0.85}
\usepackage[T1]{fontenc}
\usepackage{amsmath}
\usepackage{amsfonts}
\usepackage{amssymb}
\usepackage{array}
\usepackage{fontspec}
\usepackage{booktabs}
\usepackage{hyperref}
\usepackage[version=4]{mhchem}
\usepackage{stmaryrd}
\usepackage[dvipsnames]{xcolor}
\colorlet{LightRubineRed}{RubineRed!70}
\colorlet{Mycolor1}{green!10!orange}
\definecolor{Mycolor2}{HTML}{00F9DE}
\usepackage{graphicx}
\usepackage{amsmath}
\usepackage{graphicx}
\usepackage{capt-of}
\usepackage{lipsum}
\usepackage{fancyvrb}
\usepackage{tabularx}
\usepackage{listings}
\usepackage[export]{adjustbox}
\graphicspath{ {./images/} }
\usepackage[utf8]{inputenc}
\usepackage[english]{babel}
\usepackage{float}
\usepackage{lipsum}
\usepackage{graphicx}
\usepackage{float}
\usepackage[margin=0.7in]{geometry}
\usepackage{amsmath}
\usepackage{graphicx}
\usepackage{capt-of}
\usepackage{tcolorbox}
\usepackage{lipsum}
\usepackage{graphicx}
\usepackage{float}
\usepackage{listings}
\definecolor{deepred}{RGB}{139,0,0}
\usepackage{hyperref} 
\usepackage{xcolor} % For custom colors
\lstset{
	language=Python,                % Choose the language (e.g., Python, C, R)
	basicstyle=\ttfamily\small, % Font size and type
	keywordstyle=\color{blue},  % Keywords color
	commentstyle=\color{gray},  % Comments color
	stringstyle=\color{red},    % String color
	numbers=left,               % Line numbers
	numberstyle=\tiny\color{gray}, % Line number style
	stepnumber=1,               % Numbering step
	breaklines=true,            % Auto line break
	backgroundcolor=\color{black!5}, % Light gray background
	frame=single,               % Frame around the code
}
\usepackage{float}
\usepackage[]{amsthm} %lets us use \begin{proof}
	\usepackage[]{amssymb} %gives us the character \varnothing
	
	
	\title{Homework 2, STAT 5291}
	\author{Zongyi Liu}
	\date{Mon, Mar 2, 2026}
	\begin{document}
		\maketitle
		
		\section{Question 1}
		
		{\color{deepred}\fontspec{Georgia} Literature Review for the Final Project}
		
		
		\emph{Perform the following task}
		
		Present the following tasks in a meaningful order so that you can start including sections in your write-up. Students must review one paper for every group member. For example, if students are working in a group of three members, then students must review three papers and consequently submit three one-page write-ups. If students are working in a group of five members, then students must review five papers and consequently submit five one-page write-ups.
		
		\begin{enumerate}
			\item[i] List out your project teammates (names and UNIs).
			
			\item[ii] Find a relevant research paper related to your project topic.
			
			\item[iii] Write up a one page document summarizing the paper.
			
			\item[iv] Write a bibliography including your chosen papers.
			
			\item[v] Write an introduction section in your paper that ties in the relevant papers with your project goals and/or research question(s). This introduction can be used in your paper and will likely be edited a few times before final submission.
		\end{enumerate}
		
		\textbf{Answer}
		
		\clearpage
		
		\section{Question 2}
		{\color{deepred}\fontspec{Georgia} Wilcoxon Sum-Rank Test}
		
		\medskip
		
		Recall the Wilcoxon Sum-Rank Test Procedure:
		
		\medskip
		
		\emph{Assumptions}
		
		\medskip
		
		Assume that $\{X_1,\dots,X_m\}$ and $\{Y_1,\dots,Y_n\}$ are independent samples from different distributions with respective means $\mu_1$ and $\mu_2$. The distributions of $X$ and $Y$ have the same shape and differ only in location, i.e., by means $\mu_1$ and $\mu_2$. Note that the distributions do not have to be symmetric.
		
		\medskip
		
		\underline{Construction of the Wilcoxon Sum-Rank test statistic}
		
		\begin{enumerate}
			\item[a] Combine the two samples, and rank the observations.
			
			\item[b] Let $T_1$ be the sum of the ranks for observations in the first group, and $T_2$ be the sum of the ranks for observations in the second group.
		\end{enumerate}
		
		\medskip
		
		\emph{Perform the following tasks}
		
		\begin{enumerate}
			\item[i] Suppose that $m=n=2$. Derive the sampling distribution of $T_1$ under $H_0$. Also compute $E[T_1]$ and $Var[T_1]$ based on the derived sampling distribution.
			
			\item[ii] Suppose that $m=3$ and $n=4$. Derive the sampling distribution of $T_1$ under $H_0$. Also compute $E[T_1]$ and $Var[T_1]$ based on the derived sampling distribution.
			
			\item[iii] Now for any $m,n>0$, show that under $H_0$,
			\[
			E[T_1]=\frac{m(m+n+1)}{2}
			\]
			
			and
			
			\[
			Var[T_1]=\frac{mn(m+n+1)}{12}.
			\]
			
			\emph{Hint} Under $H_0$, the rank $R_i$ of $X_i$ is equally likely to be any one of the possible values $1,2,\dots,m+n$.
			
			\item[iv] Consider the following toy dataset which consists of a continuous response variable split by the respective groups. Equivalently, each group represents a realization from a continuous population. We are interested in the one-tailed test:
			
			\[
			H_0:\ \mu_1-\mu_2=0
			\qquad \text{versus} \qquad
			H_A:\ \mu_1-\mu_2<0.
			\]
			
			\begin{center}
				\begin{tabular}{lcccccccc}
					Group1 & 1.9168 & 3.5102 & 3.4567 & 2.7052 & 1.3366 & 2.7654 & 2.9780 & 2.8394 \\
					& 2.1455 & 1.7992 \\
					\hline
					Group2 & 2.9204 & 3.1500 & 3.3405 & 2.7374 & 4.0499 & 3.8547 & 4.1359 & 4.0159 \\
					& 3.4328 & 2.0894 & 3.7789 & 4.2853 & 2.8724 & 3.0208 & 3.1968
				\end{tabular}
			\end{center}
			
			Manually code up a function called \texttt{my\_wilcoxon()} which computes $T_1$ based on generic datasets \texttt{data\_1} and \texttt{data\_2}. Test your function on the above dataset, i.e.,
			
			\[
			\texttt{my\_wilcoxon(Group1,Group2)}.
			\]
			
			\emph{Note} To check if your function is returning the correct value, the R function
			
			\[
			\texttt{wilcox.test(Group1,Group2)}
			\]
			
			is technically computing the Mann-Whitney $U$ statistic.
			
			\item[v] Simulate the sampling distribution of $T_1$ under $H_0$ using
			
			\[
			X_1,\dots,X_m \sim f
			\qquad
			(\text{Some Continuous Distribution})
			\]
			
			and
			
			\[
			Y_1,\dots,Y_n \sim f
			\qquad
			(\text{Some Continuous Distribution})
			\]
			
			and computing each $T_1$ based on your function
			
			\[
			T_1=\texttt{my\_wilcoxon(Group1,Group2)}.
			\]
			
			You can choose any continuous distribution $f$ for Group1 or Group2 as long as they are the same under $H_0$. Simulate 10,000 draws from the sampling distribution of $T_1$ and plot a histogram of your results.
			
			
			\item[vi] Run the Wilcoxon sum-rank procedure to test the one-tailed hypothesis
			
			\[
			H_0 : \mu_1 - \mu_2 = 0 
			\qquad \text{versus} \qquad
			H_A : \mu_1 - \mu_2 < 0.
			\]
			
			Compute the p-value using the Monte Carlo sampling distribution from Part 2.v. Also compute the large sample p-value and compare the results. Note that the large sample test statistic is:
			
			\[
			z_{calc}=
			\frac{
				T_1-\dfrac{m(m+n+1)}{2}
			}{
				\sqrt{\dfrac{mn(m+n+1)}{12}}
			}.
			\]
		\end{enumerate}
		
		\textbf{Answer}
		
		\clearpage
		
		\section{Question 3}
		
		{\color{deepred}\fontspec{Georgia} Parametric and Nonparametric Two-sample Procedures}
		
		\medskip
		
		Consider the toy dataset from Problem 2, which consists of a continuous response variable split by the respective groups. Equivalently, each group represents a realization from a continuous population. We are interested in the two-tailed test:
		
		\[
		H_0 : \mu_1 - \mu_2 = 0
		\qquad \text{versus} \qquad
		H_A : \mu_1 - \mu_2 \ne 0.
		\]
		
		\begin{center}
			\begin{tabular}{lcccccccc}
				\hline
				Group1 & 1.9168 & 3.5102 & 3.4567 & 2.7052 & 1.3366 & 2.7654 & 2.9780 & 2.8394 \\
				& 2.1455 & 1.7992\\
				\hline
				Group2 & 2.9204 & 3.1500 & 3.3405 & 2.7374 & 4.0499 & 3.8547 & 4.1359 & 4.0159 \\
				& 3.4328 & 2.0894 & 3.7789 & 4.2853 & 2.8724 & 3.0208 & 3.1968
			\end{tabular}
		\end{center}
		
		\medskip
		
		\emph{Perform the following tasks}
		
		\begin{enumerate}
			
			\item[i] Construct a multiple box-plot to investigate the research question of interest. Also construct qqplots of each dataset and comment on the respective distributions.
			
			\item[ii] Run the relevant test to check if the population variances differ per group, i.e.,
			
			\[
			H_0 : \sigma_1^2 = \sigma_2^2
			\qquad \text{versus} \qquad
			H_A : \sigma_1^2 \ne \sigma_2^2.
			\]
			
			Use $\alpha = .10$ for this test.
			
			\item[iii] Based on Problem 3.ii, run a two-sample t-test or pooled two-sample t-test on the null alternative pair
			
			\[
			H_0 : \mu_1 - \mu_2 = 0
			\qquad \text{versus} \qquad
			H_A : \mu_1 - \mu_2 \ne 0.
			\]
			
			Use $\alpha = .05$ for this test. Make sure to include all relevant steps in the testing procedure for full credit.
			
			\item[iv] Run a bootstrap procedure to test the hypothesis
			
			\[
			H_0 : \mu_1 - \mu_2 = 0
			\qquad \text{versus} \qquad
			H_A : \mu_1 - \mu_2 \ne 0.
			\]
			
			Make sure to include all relevant steps in the testing procedure for full credit.
			
			\item[v] Run a permutation test procedure to test the hypothesis
			
			\[
			H_0 : \mu_1 - \mu_2 = 0
			\qquad \text{versus} \qquad
			H_A : \mu_1 - \mu_2 \ne 0.
			\]
			
			Make sure to include all relevant steps in the testing procedure for full credit.
			
		\end{enumerate}
		
		\textbf{Answer}
		
		
		\clearpage
		
		\section{Question 4}
		
		{\color{deepred}\fontspec{Georgia} Power and Sample Size Determination}
		
		\medskip
		
		Consider the two-sample pooled t-procedure for testing the upper tailed hypothesis:
		
		\[
		H_0 : \mu_1 - \mu_2 = \Delta_0
		\qquad \text{versus} \qquad
		H_A : \mu_1 - \mu_2 > \Delta_0.
		\]
		
		\medskip
		
		\emph{Perform the following tasks}
		
		\begin{enumerate}
			
			\item[i] Derive an expression for the power of the upper tailed test as a function of the non-centrality parameter $\delta$. The power derivation should be similar to the class example for the one-sample t-test. Students should also identify the correct expression for $\delta$.
			
			\item[ii] Suppose that the experimenter would like to determine the minimal sample size required of the pooled two-sample t-test in order to achieve a pre-specified power $1-\beta$. More specifically, the experimenter uses a pilot study to estimate the effect size
			
			\[
			d=\frac{\mu_2-\mu_1}{\sigma},
			\]
			
			and the non-centrality parameter $\delta$ from Problem 4.i. For simplicity, we also assume a balanced design $(m=n)$.
			
		\end{enumerate}
		
		\medskip
		
		\emph{Required task}
		
		\medskip
		
		For fixed $\alpha$ and fixed effect size $d$ (or fixed non-centrality parameter $\delta$), compute the power curve as a function of the sample size $n$, i.e., construct the curve
		
		\[
		\text{power}(n), \qquad \text{for } n=2,3,\dots,200,
		\]
		
		where $\alpha$ and $d$ are fixed.
		
		Students are required to compute power curves related to three different pilot studies provided in the files \texttt{HW2\_P4\_data1.csv}, \texttt{HW2\_P4\_data2.csv} and \texttt{HW2\_P4\_data3.csv}. Further, students should compute the power curves for two significance levels $\alpha=.05$ and $\alpha=.10$. Your final answer should be summarized in one elegant plot showing six power curves. Please label the legend appropriately to clearly identify which curve corresponds to the $d$ (or $\delta$) and $\alpha$.
		
		\medskip
		
		\emph{Note} maybe write a loop in R or Python to compute these curves?
		
		\textbf{Answer}
		
		\clearpage
		
		\section{Question 5}
		{\color{deepred}\fontspec{Georgia} Two-Sample Pooled Test and Likelihood Ratio}
		
		Recall the pooled two-sample t-procedure, which assumes the model:
		
		\medskip
		
		\underline{Assumptions}
		
		\begin{enumerate}
			\item $X_1, X_2, \ldots, X_m$ is a random sample from a normal distribution with mean $\mu_1$ and variance $\sigma^2$. The variance $\sigma^2$ is unknown.
			
			\item $Y_1, Y_2, \ldots, Y_n$ is a random sample from a normal distribution with mean $\mu_2$ and variance $\sigma^2$. The variance $\sigma^2$ is unknown.
			
			\item The $X$ and $Y$ samples are independent of one another.
		\end{enumerate}
		
		\medskip
		
		Consider the two tailed hypothesis
		\[
		H_0 : \mu_1 - \mu_2 = 0 
		\qquad \text{versus} \qquad
		H_A : \mu_1 - \mu_2 \neq 0 .
		\]
		
		\medskip
		
		\emph{Perform the following tasks}
		
		\begin{enumerate}
			
			\item[i] Derive an expression for the likelihood
			\[
			\mathcal{L}(\theta) = \mathcal{L}(\mu_1,\mu_2,\sigma^2).
			\]
			
			\item[ii] Find the maximum likelihood estimator of
			\[
			\theta = (\mu_1,\mu_2,\sigma^2).
			\]
			
			Note that you are not required to perform the \emph{second derivative test} for this problem.
			
			\item[iii] Derive the likelihood ratio test statistic $\lambda$. Note that you are not required to perform the \emph{second derivative test} for this problem.
			
			\item[iv] Show that the likelihood ratio test can be expressed in terms of the t-statistic $(T)$, i.e., there exists a function $f$ such that
			\[
			\lambda = f(T).
			\]
			
		\end{enumerate}
		
		\textbf{Answer}
		
		
		\clearpage
		
		
		\section{Question 6}
		
		{\color{deepred}\fontspec{Georgia} Parametric and Non-parametric Paired Procedures}
		
		\medskip
		
		Consider testing the stopping distance [ft] of two types of tires on 6 different cars. 
		The response variable is the stopping distance of the car after the brakes are applied 
		at 60 mph. The two types of tires are: Winter and High Performance. 
		Does the data set indicate that the stopping distance differs between the two tire types?
		
		\medskip
		
		\begin{center}
			\begin{tabular}{ccc}
				\hline
				Car & High Performance & Winter \\
				\hline
				1 & 226.6 & 239.5 \\
				2 & 230.5 & 257.5 \\
				3 & 220.0 & 236.8 \\
				4 & 215.9 & 223.9 \\
				5 & 231.5 & 248.0 \\
				6 & 235.5 & 253.6 \\
				\hline
			\end{tabular}
		\end{center}
		
		\medskip
		
		\emph{Perform the following tasks}
		
		\begin{enumerate}
			
			\item[i] Assuming the stopping distance for both tire types are approximately normally distributed, run the paired two-sample t-procedure to test the claim
			
			\[
			H_0 : \mu_D = 0 
			\qquad \text{versus} \qquad
			H_A : \mu_D \neq 0 .
			\]
			
			Make sure to include all relevant steps in the testing procedure for full credit.
			
			\item[ii] Assuming the stopping distance for $D = X - Y$ is symmetric, run the Wilcoxon signed-rank procedure to test the claim
			
			\[
			H_0 : \mu_D = 0 
			\qquad \text{versus} \qquad
			H_A : \mu_D \neq 0 .
			\]
			
			Make sure to include all relevant steps in the testing procedure for full credit.
			
		\end{enumerate}
		
		
		\textbf{Answer}
		
		
		\clearpage
		
		\section{Question 7}
		
		{\color{deepred}\fontspec{Georgia} Categorical Hypothesis Testing Problem}
		
		\medskip
		
		Binomial data gathered from more than one population are often presented in a contingency table. 
		For the case of two populations, the table might look like
		
		\begin{center}
			\begin{tabular}{c|c|c|c}
				& Population 1 & Population 2 & Total \\
				\hline
				Successes & $S_1$ & $S_2$ & $S = S_1 + S_2$ \\
				Failures & $F_1$ & $F_2$ & $F = F_1 + F_2$ \\
				\hline
				Total & $m$ & $n$ & $N = m + n$
			\end{tabular}
		\end{center}
		
		where Population 1 is $\text{Binomial}(m,p_1)$, with $S_1$ successes and $F_1$ failures, and Population 2 is $\text{Binomial}(n,p_2)$, with $S_2$ successes and $F_2$ failures. 
		We will motivate both the two-sample $z$-test for proportions and the chi-squared test for association based on the $2\times2$ contingency table. 
		The hypothesis of interest is:
		
		\[
		H_0: p_1 = p_2
		\qquad \text{versus} \qquad
		H_A: p_1 \neq p_2
		\]
		
		\emph{Perform the following tasks}
		
		\begin{enumerate}
			
			\item[i] Show that a test can be based on the statistic
			\[
			W = \frac{(\hat p_1 - \hat p_2)^2}{\left(\frac{1}{m} + \frac{1}{n}\right)\left(\hat p (1-\hat p)\right)},
			\]
			where $\hat p_1 = S_1/m$, $\hat p_2 = S_2/n$, and
			\[
			\hat p = \frac{S_1 + S_2}{m+n}.
			\]
			Also show that as $m,n \to \infty$, the distribution of $W$ approaches $\chi^2_1$.
			
			\item[ii] Another way of measuring the departure from $H_0$ is by calculating the expected frequency table. 
			This table is constructed by conditioning on the marginal totals and filling in the table according to $H_0: p_1=p_2$, that is,
			
			\medskip
			
			\begin{center}
				Expected Frequencies
				
				\begin{tabular}{c|c|c|c}
					& Population 1 & Population 2 & Total \\
					\hline
					Successes & $\frac{mS}{m+n}$ & $\frac{nS}{m+n}$ & $S = S_1 + S_2$ \\
					Failures & $\frac{mF}{m+n}$ & $\frac{nF}{m+n}$ & $F = F_1 + F_2$ \\
					\hline
					Total & $m$ & $n$ & $N = m+n$
				\end{tabular}
			\end{center}
			
			Using the expected frequency table, a statistic $W^*$ is computed by going through the cells of the tables and computing
			
			\[
			W^* =
			\sum \frac{(\text{observed} - \text{expected})^2}{\text{expected}}
			\]
			
			\[
			=
			\frac{\left(S_1 - \frac{mS}{m+n}\right)^2}{\frac{mS}{m+n}}
			+
			\frac{\left(S_2 - \frac{nS}{m+n}\right)^2}{\frac{nS}{m+n}}
			+
			\frac{\left(F_1 - \frac{mF}{m+n}\right)^2}{\frac{mF}{m+n}}
			+
			\frac{\left(F_2 - \frac{nF}{m+n}\right)^2}{\frac{nF}{m+n}}.
			\]
			
			Show, algebraically, that $W = W^*$ and hence that $W^*$ is asymptotically chi-squared.
			
			\item[iii] A famous medical experiment was conducted by Joseph Lister in the late 1800s. 
			Mortality associated with surgery was quite high, and Lister conjectured that the use of disinfectant, carbolic acid, would help. 
			Over a period of several years Lister performed 75 amputations with and without carbonic acid. 
			The data are given here:
			
			\begin{center}
				Carbolic acid Used?
				
				\begin{tabular}{c|c|c}
					& Yes & No \\
					\hline
					Patient lived & 34 & 19 \\
					Patient died & 6 & 16
				\end{tabular}
			\end{center}
			
			Use these data to test whether the use of carbolic acid is associated with patient mortality.
			
		\end{enumerate}
		
		
		\textbf{Answer}
		
		
		\clearpage
		
		
		\section{Question 8}
		
		{\color{deepred}\fontspec{Georgia} Goodness-of-fit (Practice Problem)}
		
		\medskip
		
		Construct a Chi-squared goodness-of-fit test to assess whether or not a dataset $(X)$ is normally distributed, i.e., test the null alternative pair:
		
		\[
		H_0 : X \sim \text{Normal}
		\qquad \text{versus} \qquad
		H_A : X \text{ is not Normal}.
		\]
		
		To accomplish this task, you must first convert the continuous observed data into a frequency table, i.e., the heights of the bars on a histogram. 
		Then apply the goodness-of-fit test based on the observed and expected counts assuming the null hypothesis $H_0: X\sim\text{Normal}$ is true.
		
		Some considerations follow:
		
		\begin{itemize}
			
			\item Let $K$ be the number of \emph{bins} in the histogram. Note that changing the number of bins will consequently change the degrees of freedom of the chi-squared distribution.
			
			\item Don't forget that you have to estimate both $\mu$ and $\sigma^2$ in order to compute the expected counts, which also impacts the degrees of freedom of the chi-squared distribution.
			
			\item If you are successful in solving this problem, maybe simulate a few normal and non-normal datasets and compare the qqplots with the goodness-of-fit outcomes of the simulated data.
			
		\end{itemize}
		
		\textbf{Answer}
		
		
		\clearpage
		
	\end{document}hiont
